% paper/sections/13_conclusion.tex
\section{まとめと今後の課題}
\label{sec:conclusion}

% ═══════════════════════════════════════════════════════════════════════════════
\subsection{本稿の貢献のまとめ}
\label{sec:contribution_summary}
% ═══════════════════════════════════════════════════════════════════════════════

本稿では，気液二相流の高精度数値シミュレーションに必要な数学的・数値的基盤を，
学部生・大学院生を主な読者として体系的に解説した．
本稿の技術的中核は\textbf{寄生流れの根本的抑制}にあり，
CCD $O(h^6)$ による高精度曲率計算，
FD 直接法 PPE + CCD 勾配による Balanced-Force 条件の $O(h^6)$ 維持，
および DCCD 散逸機構がその主要手段である（§\ref{sec:intro}）．
各章の設計方針を整理する前に，「なぜこの組み合わせが必要か」を再確認する．

\begin{center}
\renewcommand{\arraystretch}{1.4}
\begin{tabular}{>{\raggedright}p{60mm}>{\raggedright\arraybackslash}p{70mm}}
\hline
\textbf{数値的困難} & \textbf{本稿の対処手法} \\
\hline
界面の質量損失       & CLS 法（第\ref{sec:levelset}章）\\
低次移流による界面にじみ & Dissipative CCD + TVD-RK3（第\ref{sec:advection}章）\\
曲率計算誤差 $\to$ 寄生流れ & CCD 法 $O(h^6)$（第\ref{sec:CCD}章）\\
チェッカーボード圧力振動 & DCCD 散逸（$\varepsilon_d=1/4$，§\ref{sec:dccd_decoupling}，第\ref{sec:collocate}章）\\
変密度 PPE の収束悪化 & 分相 PPE 解法（§\ref{sec:field_extension}，§\ref{sec:varrho_ppe_limitation}）\\
格子粗さによる界面解像不足 & Level Set 連動非一様格子（第\ref{sec:grid}章）\\
\hline
\end{tabular}
\end{center}

% ─────────────────────────────────────────────────────────────────────────────
\subsubsection{各章の技術的貢献}
% ─────────────────────────────────────────────────────────────────────────────

\begin{description}
  \item[\textbf{第\ref{sec:governing}章：支配方程式と One-Fluid 定式化}]
        2つの独立した Navier--Stokes 方程式系（液相・気相）を，
        ヘヴィサイド関数 $\psi$ を用いた一体化方程式に変換する
        One-Fluid 定式化を導出した（§\ref{sec:onefluid}）．
        界面ジャンプ条件 $[p]_\Gamma=\kappa/We$（Young--Laplace 則）の離散的取り扱いを
        CSF 体積力形式（付録~\ref{app:csf_full}）とHFE（§\ref{sec:field_extension}）の
        両面から整理した．
        無次元化（Re，Fr，We の定義）と曲率公式の幾何学的導出（§\ref{sec:curvature}）により，
        後の章の数値実装の数学的基礎を確立した．

  \item[\textbf{第\ref{sec:levelset}章：Conservative Level Set 法}]
        標準 Level Set 法の質量損失問題を克服するため，
        $\psi = H_\varepsilon(\phi) \in [0,1]$ を移流変数とする
        \textbf{保存形}の界面追跡法（CLS 法）を定式化した．
        再初期化方程式（圧縮--拡散 PDE）により，
        移流で崩れた界面プロファイルを $H_\varepsilon(\phi)$ の形に回復する理論を示した．
        $\psi \to \phi$ の解析的逆変換（§\ref{sec:newton_inversion}）により，
        幾何計算に必要な距離関数 $\phi$（$|\nabla\phi|=1$）の高精度復元手順を明示した．

  \item[\textbf{第\ref{sec:advection}章：CLS 移流の離散化（Dissipative CCD スキーム・時間積分・安定性制約）}]
        純粋な CCD 中心差分移流のフォン・ノイマン不安定性（$|g(\xi)|^2 > 1$）を証明し，
        スペクトルフィルタ $H(\xi;\varepsilon_d) = 1 - 4\varepsilon_d\sin^2(\xi/2)$ による
        振幅抑制機構を定式化した（§\ref{sec:advection_motivation}）．
        CLS 移流には一様 $\varepsilon_d^{\text{adv}} = 0.05$ を採用し，
        界面での Gibbs 振動を効果的に抑制する設計根拠を示した．
        CLS 移流・再初期化には TVD-RK3（Shu--Osher 1988）を適用した．
        AB2 対流外挿と IPC スプリッティングの組み合わせにより
        \textbf{全体時間精度 $O(\Delta t^2)$} を達成したことを明示した（§\ref{sec:time_int}）．
        毛管波制約 $\Delta t_\sigma \propto \Delta x^{3/2}$ を分散関係から導出した（§\ref{sec:stability}）．
        WENO5 スキームの理論と実装詳細は付録~\ref{app:weno5} に参照スキームとして収録した．

  \item[\textbf{第\ref{sec:CCD}章：結合コンパクト差分法（CCD 法）}]
        3点ステンシルで1階・2階微分を同時求解する CCD 法（Chu \& Fan 1998）を定式化し，
        係数（$\alpha_1 = 7/16$，$a_1 = 15/16$，$b_1 = 1/16$ 等）の導出根拠を示した
        （詳細な Taylor 展開証明は付録~\ref{app:ccd_coef_derivation}）．
        CCD 法を楕円型境界値問題のソルバーとして用いる
        \textbf{仮想時間陰解法}（§\ref{sec:ccd_elliptic}）を定式化し，
        変密度 PPE への応用（§\ref{sec:ccd_poisson_matrix}）を示した．

  \item[\textbf{第\ref{sec:grid}章：界面適合非一様格子}]
        Level Set 値に連動したグリッド密度関数 $\omega(\phi)$ により，
        界面近傍のみを自動的に高解像度化する非一様テンソル積格子を定式化した（§\ref{sec:grid_gen}）．
        座標変換メトリック係数の CCD 法による $O(h^6)$ 精度評価と
        ALE 効果の簡略化条件を明示した（§\ref{sec:grid_ale}）．

  \item[\textbf{第\ref{sec:collocate}章：圧力--速度連成と DCCD 散逸}]
        コロケート格子でのチェッカーボード不安定のメカニズム（$2\Delta x$ モード解析）を示し，
        Dissipative CCD 散逸（$\varepsilon_d=1/4$，§\ref{sec:dccd_decoupling}）によるチェッカーボード
        抑制の設計原理を示した．Rhie--Chow 補正（付録~\ref{app:rhie_chow_full}）および
        Balanced-Force 条件（§\ref{sec:balanced_force}）の理論的根拠を参照収録した．
        さらに，CCD が $p'$ と同時に返す $p''$ を用いた係数相殺補正（§~\ref{sec:rc_high_order}）により，
        RC ブラケットの精度を $\Ord{h^2}$ から $\Ord{h^4}$ に\textbf{追加コストゼロ}で改善できることを示し，
        格子収束テスト（§~\ref{sec:verify_rc_bracket}）で数値的に確認した．

  \item[\textbf{第\ref{sec:pressure}章：変密度圧力 Poisson 方程式}]
        変密度 IPC Projection 法から
        変係数 PPE $\nabla\cdot(1/\rho\,\nabla p) = (1/\Delta t)\nabla\cdot\bu^*$ を導出し，
        FD 直接法による PPE 求解と CCD 勾配の併用
        （§\ref{sec:ppe_discretization_choice}）を本稿の標準求解手法として採用した：
        PPE を 2 次精度 FD で厳密に解き（残差ゼロ），
        速度補正の勾配に CCD $\Ord{h^6}$ を適用することで
        Balanced-Force 条件を $\Ord{h^6}$ 精度で維持する．
        CCD PPE の欠陥補正法（§\ref{sec:defect_correction_main}）も定式化したが，
        $L_L^{-1} L_H$ の不定性により定常反復法では収束不能であり
        （§\ref{sec:ppe_discretization_choice}），
        将来の Krylov 法前処理との組み合わせに向けた基盤技術として位置づける．
        HFE（§\ref{sec:field_extension}）は分相 PPE 解法の基盤技術として定式化・検証した．

  \item[\textbf{第\ref{sec:algorithm}章：完全アルゴリズムフロー}]
        NS 方程式の各項と離散ソルバーの対応関係（§\ref{sec:algo_operator_map}），
        7ステップの全体フロー（§\ref{sec:algo_flow}），
        および CFL・毛管波制約を統合したタイムステップ制御（§\ref{sec:algo_timestep}）を整理した．
        AB2 + IPC による対流項処理・Crank--Nicolson 粘性項の ADI 求解・
        IPC 増分圧力 Poisson 求解の完全な離散形式を §\ref{sec:algorithm} に示した．
\end{description}

% ─────────────────────────────────────────────────────────────────────────────
\subsubsection{精度バランスの整理}
% ─────────────────────────────────────────────────────────────────────────────

本稿の数値手法における各コンポーネントの精度は，
第\ref{sec:pressure}章の表\ref{tab:accuracy_summary}に整理した通りである．
ここでは，その設計上のトレードオフを技術的貢献の観点から振り返る．
主要な精度設計上のトレードオフとして，
（1）空間精度は CCD $\Ord{h^6}$・Dissipative CCD $O(\varepsilon_d h^2)$（フィルタ律速）と高次で統一され，
PPE は FD 直接法 $\Ord{h^2}$ + CCD 勾配 $\Ord{h^6}$ の併用で
Balanced-Force 条件を $\Ord{h^6}$ 精度で維持するが，
CSF $O(\varepsilon^2)$ モデル誤差の根本的除去には分相 PPE 解法（§\ref{sec:future_work}）が必要である（付録~\ref{app:csf_full}）；
（2）時間精度は AB2+IPC スキームにより $O(\Delta t^2)$ を達成しており，
さらなる高次化には $O(\Delta t^3)$ 以上の高次 Projection 法（例：Guermond--Shen 2006）への移行が必要である，
の2点が挙げられる（詳細は§\ref{sec:pressure_accuracy_summary}参照）．
これらの特性は意図的な設計選択であり，
空間全コンポーネントの高次精度統一と曲率計算の高精度化が寄生流れを直接抑制するという工学的意義に基づく．

% ═══════════════════════════════════════════════════════════════════════════════
\subsection{今後の課題}
\label{sec:future_work}
% ═══════════════════════════════════════════════════════════════════════════════

本稿の定式化（2次元・等温・非圧縮・固定格子）の主要な拡張方向を示す．

% ─────────────────────────────────────────────────────────────────────────────
\subsubsection{数値精度の向上}
% ─────────────────────────────────────────────────────────────────────────────

\begin{description}
  \item[\textbf{NS 方程式の時間高精度化：}]
        本稿では AB2+IPC スキーム（van Kan 1986）により全体時間精度 $O(\Delta t^2)$ を達成している
        （§\ref{sec:time_accuracy_table}）．
        さらに $O(\Delta t^3)$ 以上の時間精度へ向上させるには，
        スプリッティング誤差が $O(\Delta t^3)$ 以上の高次 Projection 法への移行が必要である．
        具体的には Kim--Moin 法（1985）の3次版や Guermond--Shen 法（2006）が候補であり，
        各法の安定性・実装コストと精度向上の比較が今後の研究対象となる．

  \item[\textbf{PPE 求解精度の更なる検証：}]
        本稿では CCD-PPE（$\Ord{h^6}$）の等密度 Poisson での格子収束を検証した（§\ref{sec:verify_ppe}）．
        今後の課題として，高密度比（$\rho_l/\rho_g \geq 10^3$）・粗格子での収束挙動の詳細解析，
        および並列計算環境での CCD トリダイアゴナルソルバーのスケーラビリティ評価が挙げられる．

  \item[\textbf{表面張力の陰的処理：}]
        毛管波制約 $\Delta t_\sigma \propto \Delta x^{3/2}$ を除去するため，
        Implicit Surface Tension 法（例：Raessi \& Pitsch 2012）の導入が有効である．
        表面張力項を陰的に扱うことで，時間刻み制約を対流 CFL 条件のみとしながら
        安定計算が可能となり，高解像度計算での計算コストを大幅に削減できる．

  \item[\textbf{分相 PPE 解法による変密度 PPE の収束性回復：}]
        分相 PPE 解法の導入は，精度向上と変密度 PPE 収束性の双方の観点から\textbf{最優先課題}である．

        \textbf{精度面：}
        §\ref{sec:pressure_accuracy_summary} で整理した通り，
        CSF モデルの近似誤差 $\Ord{\varepsilon^2} \approx \Ord{\Delta x^2}$ が
        全空間コンポーネント（CCD-PPE $\Ord{h^6}$，CCD 移流 $\Ord{h^6}$）の律速段階である．
        Kang ら~\cite{Kang2000} の Ghost Fluid Method に倣い，
        圧力ジャンプ条件 $[p] = \sigma\kappa$ を PPE 離散化式にゴースト値として直接組み込むことで，
        CSF の界面厚み $\varepsilon$ に起因する $\Ord{\Delta x^2}$ モデル誤差を原理的に排除できる．

        \textbf{収束性：}
        §\ref{sec:varrho_ppe_limitation} で示した通り，
        界面型密度場（smoothed Heaviside，$\rho_l/\rho_g \geq 10$）では
        変密度 PPE の欠陥補正法が発散する．
        根本原因は $1/\rho$ の界面ジャンプを横切る FD/CCD ステンシルの $\Ord{1}$ 離散化誤差である．
        各相で定数密度 PPE（$\rho_k^{-1}\nabla^2 p = q$）を解き，
        界面のジャンプ条件を HFE（§\ref{sec:field_extension}）で課すことで，
        不連続係数を横切るステンシルが排除され，
        各相の CCD 離散化が $\Ord{h^6}$ の設計精度を回復する．

  \item[\textbf{DCCD 内在的散逸による Rhie--Chow 補正の排除：}]
        Rhie--Chow 補正は $2\Delta x$ チェッカーボードモードを抑制する標準的手法であるが，
        補正項が界面近傍の Balanced-Force 条件（§\ref{sec:balanced_force}）を
        局所的に乱し，寄生流れの一因となる副作用がある．
        DCCD の散逸パラメータ $\beta$ を波数 $k\Delta x = \pi$ において
        振幅増幅率が $0$ に収束するよう設計すれば，
        DCCD 自身がデカップリングモードを自然に抑制するため，
        外部補正項としての Rhie--Chow 項が不要となる．
        さらに，SAT（Simultaneous Approximation Term）法により
        境界節点での数値エネルギー流入をペナルティ項として制御することで，
        全計算領域にわたる安定なデカップリング抑制が保証される．
        この方針は本稿の Balanced-Force 設計との整合性が高く，
        寄生流れのさらなる低減が期待される．
\end{description}

% ─────────────────────────────────────────────────────────────────────────────
\subsubsection{機能拡張}
% ─────────────────────────────────────────────────────────────────────────────

\begin{description}
  \item[\textbf{3次元実装と並列化：}]
        CCD 演算子・Rhie--Chow 補正・PPE 求解の3方向（$x,y,z$）拡張は
        概念的には直接的であるが，実装では MPI による非局所通信（ハロー交換）が必要となる．
        3D CCD の LU 分解コストと並列スケーラビリティの評価が課題である．

  \item[\textbf{適応格子細分化（AMR）との統合：}]
        界面近傍の動的細分化（AMR）は計算コストの大幅削減をもたらす．
        粗・密格子境界での CCD ステンシルの整合性，
        および非一様格子上での Dissipative CCD フィルタ適用が技術的課題となる．

  \item[\textbf{相変化の組み込み：}]
        沸騰・蒸発問題への対応には，
        界面でのエネルギーバランス（Stefan 条件）とエネルギー方程式の連成が必要である．
        界面での質量フラックスが発散条件（$\nabla\cdot\bu = 0$）を修正するため，
        Projection 法の定式化も修正が必要となる．

  \item[\textbf{乱流モデルとの統合：}]
        高 Reynolds 数の気液二相乱流では，界面近傍のサブグリッドスケールモデリングが重要である．
        LES/RANS との結合では，界面近傍での渦粘性モデルの適用範囲
        （「界面の内側か外側か」）の定義が主要な研究課題となる．

  \item[\textbf{非等温・非平衡問題：}]
        本稿は等温を仮定しているが，凝縮・蒸発・燃焼などの非等温現象では
        温度場・化学種輸送方程式との連成が必要となる．
        物性値の温度依存性（$\sigma(T)$，$\mu(T)$）も
        Marangoni 流れ（表面張力勾配駆動流）などの新たな物理を生む．
\end{description}

% ─────────────────────────────────────────────────────────────────────────────
\subsubsection{検証と実証}
% ─────────────────────────────────────────────────────────────────────────────

\begin{description}
  \item[\textbf{多相流ベンチマークの数値実施と定量的比較：}]
        コンポーネント検証（第\ref{sec:verification}章：CCD 精度・GCL・Zalesak・PPE）は設計通りに実施済みであるが，
        物理ベンチマーク（第\ref{sec:validation}章：静止液滴・毛細管波・気泡上昇・Rayleigh--Taylor）については，
        ベンチマーク設定と目標値を第\ref{sec:validation}章に詳述した段階であり，
        フル NS ソルバーによる数値実験の実施と，
        解析解・文献値（Hysing ら \cite{Hysing2009}，Prosperetti \cite{Prosperetti1981} 等）との
        定量的比較が最優先の検証課題である．
        これにより，本稿の設計方針（Balanced-Force + CCD-PPE）の
        有効性を数値的に実証できる．

  \item[\textbf{格子収束解析による理論精度の実証：}]
        各ベンチマークについて $N = 32, 64, 128, 256$ での系統的な収束解析を行い，
        曲率 $O(h^6)$・Dissipative CCD 移流 $O(\varepsilon_d h^2)$（フィルタ律速）・寄生流れ $O(h^2)$（CSF 律速）という
        理論精度との整合を定量的に確認する．
\end{description}

% ═══════════════════════════════════════════════════════════════════════════════
\subsection{学習者へのメッセージ}
\label{sec:message}
% ═══════════════════════════════════════════════════════════════════════════════

本稿を最後まで読み通した読者は，気液二相流 CFD の主要な数値的困難と
その克服手法を体系的に習得したことになる．
しかし，本稿は「スタート地点」に過ぎない．

気液二相流の数値計算は，界面の物理（表面張力，接触角，相変化），
数値スキームの精度と安定性，そして高性能計算の工学的要請が
複雑に絡み合う，活発な研究フロンティアである．
本稿で学んだ基礎—One-Fluid 定式化，CLS 法，CCD 法，Projection 法—は，
いずれも最新の研究論文にも引き続き登場する普遍的な技術要素である．

次のステップとして，以下を推奨する：
\begin{enumerate}
  \item 本稿の手法を一つ実装し，§\ref{sec:verification}（コンポーネント検証）および §\ref{sec:validation}（多相流ベンチマーク）で段階的に検証する
  \item Tryggvason \textit{et al.} \cite{Tryggvason2011} などの
        二相流 CFD の標準テキストで知識を補完する
  \item 文献 \cite{OlssonKreiss2005}（CLS），\cite{JiangShu1996}（WENO5），
        \cite{ChuFan1998}（CCD）の原論文を熟読し，省略された数学的詳細を習得する
\end{enumerate}

本稿の定式化が，気液二相流数値解析の研究・教育に貢献することを期待する．
